
\begin{abstract}
This work develops a detailed analytical and geometric framework for the
quadratic integer sequence
\[
k(n)=\tfrac{1}{6}(2n^2+3n+1),
\]
which takes integer values precisely for the congruence classes
\(n \equiv 1,5 \pmod{6}\).

A fine-grained analysis of the first differences reveals an alternating pattern
of \emph{minor} and \emph{major steps} whose pairwise sums---as proved in this
work---form the minimal even block size for which the difference sequence yields an
arithmetic progression of circle circumferences.
This minimality, combined with the circumference interpretation, determines a constant radial increment
$\Delta R = 12/\pi$ and leads, with no remaining free parameter once one revolution is assigned to each pair, to an
explicit closed parametrisation of an Archimedean spiral with slope
$a = 6/\pi^2 = 1/\zeta(2)$, whose uniform radial growth is obtained directly
from the discrete arithmetic data of the sequence.

Extending the planar spiral by a linear height function produces a
three-dimensional conical helix on a right circular cone with half-angle
$\alpha = \arctan(2/\pi) \approx 32.48^\circ$. The algebraic core of the
entire construction is the fundamental identity
$(4n+3)^2 = 48\,k(n)+1$, which defines a universal quantity
$Q(k) = \sqrt{48k+1}$ from which radius, azimuth, and height of the
helix can be expressed as exact closed functions of a single parameter.

As a corollary of the congruence-class structure, all primes $p>3$---which by
a classical result of number theory satisfy $p \equiv 1,5 \pmod{6}$---appear
as a distinguished subset of the integer-indexed points on the helix. The helix
furthermore contains all composite numbers of these congruence classes; the
primes form a infinite subset of relative density zero therein --- one that the geometry does not separate from the composites.

This investigation demonstrates how a simple quadratic form generates a coherent
geometric representation encompassing Archimedean spirals, conical helices,
uniform discretisation, and closed analytical expressions. It illustrates how
discrete modular patterns can be embedded into continuous geometric structures,
thereby opening a unified perspective on the interplay between integer sequences,
modular arithmetic, and parametric geometry.
The contribution lies not in new number theory but in a geometric synthesis: known arithmetic building blocks are connected along a single derivation path.
\end{abstract}

\vspace{1em}
\noindent\textbf{MSC 2020:} 11B25 (Arithmetic progressions), 53A04 (Curves in Euclidean
and related spaces), 11A41 (Primes), 11B83 (Special sequences and polynomials).

\noindent\textbf{Keywords:} Quadratic integer sequence, Archimedean spiral,
conical helix, congruence classes, prime numbers, parametric geometry,
arc-length parametrisation, Basel problem $\zeta(2)$.

\clearpage
